\documentclass[12pt,a4paper]{article}
\usepackage[utf8]{vietnam}
\usepackage{amsmath,amssymb}
\usepackage{geometry}
\geometry{
  top=12mm,
  bottom=12mm,
  left=14mm,
  right=14mm,
  footskip=6mm,
  headheight=14pt
}

\usepackage{tikz}
\usetikzlibrary{arrows.meta,calc,shapes.geometric}
\usepackage[most]{tcolorbox}
\usepackage{fontawesome}
\usepackage{fancyhdr}
\usepackage{enumitem}

\pagestyle{fancy}
\fancyhf{}
\fancyhead[L]{\footnotesize\color{teal!80!black}\faBook\ \textbf{SỔ TAY KIẾN THỨC TOÁN 3}}
\fancyhead[R]{\footnotesize\color{teal!80!black}\textbf{Lớp Toán Thầy Lĩnh -- loptoanthaylinh.vn}}
\fancyfoot[C]{\footnotesize\color{gray}-- Sổ tay kiến thức: Bài 24. Gấp một số lên nhiều lần --}
\renewcommand{\headrulewidth}{0.4pt}
\renewcommand{\footrulewidth}{0pt}

\setlength{\parindent}{0pt}
\setlength{\parskip}{0pt}

\begin{document}

% --- BANNER BÀI HỌC ---
\begin{tcolorbox}[
  enhanced,
  colback=teal!8!white,
  colframe=teal!70!black,
  arc=3mm, boxrule=1.2pt,
  left=3mm, right=3mm, top=2mm, bottom=2mm,
  halign=center
]
  {\LARGE\bfseries\color{teal!90!black}\faBookmark\ BÀI 24: GẤP MỘT SỐ LÊN NHIỀU LẦN}\\[1mm]
  {\itshape\footnotesize\color{black!75} Kĩ năng vận dụng phép nhân vào bài toán thực tế, phân biệt đại lượng gấp lần trên 1 trang A4}
\end{tcolorbox}

\vspace{-2mm}

\noindent
\begin{minipage}[t]{0.49\textwidth}\vspace{0pt}
\begin{tcolorbox}[
  enhanced,
  colback=blue!4!white,
  colframe=blue!75!black,
  arc=2mm, boxrule=0.8pt,
  left=2mm, right=2mm, top=1.8mm, bottom=1.8mm,
  title={\footnotesize\bfseries\color{white}\faMagic\ 1. Quy tắc cốt lõi}
]
{\footnotesize
\begin{itemize}\setlength{\itemsep}{1pt}\setlength{\leftskip}{-2.5mm}
  \item Muốn gấp số $A$ lên $k$ lần, ta thực hiện:
  $$\text{Số mới} = A \times k$$
  \item \textbf{Bản chất:} Số mới bằng $k$ lần số ban đầu gộp lại.
  \item Đơn vị của kết quả là \textbf{đơn vị của đại lượng ban đầu $A$}, không phải là chữ "lần".
  \item \textit{Ví dụ:} Đoạn thẳng dài $12\text{ cm}$, gấp lên $3$ lần thì dài $12 \times 3 = \mathbf{36\text{ cm}}$.
\end{itemize}
}
\end{tcolorbox}
\end{minipage}
\hfill
\begin{minipage}[t]{0.49\textwidth}\vspace{0pt}
\begin{tcolorbox}[
  enhanced,
  colback=teal!4!white,
  colframe=teal!75!black,
  arc=2mm, boxrule=0.8pt,
  left=2mm, right=2mm, top=1.8mm, bottom=1.8mm,
  title={\footnotesize\bfseries\color{white}\faSliders\ 2. Kĩ thuật vẽ sơ đồ đoạn thẳng}
]
{\footnotesize
\begin{itemize}\setlength{\itemsep}{1pt}\setlength{\leftskip}{-2.5mm}
  \item \textbf{Quy ước vẽ sơ đồ chuẩn:}
  \begin{itemize}\setlength{\itemsep}{0.5pt}\setlength{\leftskip}{-2.5mm}
    \item Số bé (đại lượng ban đầu): Vẽ $1$ đoạn thẳng mẫu.
    \item Số lớn (đại lượng gấp lên): Vẽ đúng $k$ đoạn thẳng bằng nhau như thế ghép nối tiếp.
  \end{itemize}
  \item Đánh dấu ngoặc nhọn ghi tổng hoặc hiệu hai đại lượng nếu bài toán yêu cầu.
\end{itemize}
}
\end{tcolorbox}
\end{minipage}

\vspace{-2.5mm}
% --- KHỐI 3: SƠ ĐỒ TƯ DUY TIKZ ---
\noindent
\begin{tcolorbox}[
  enhanced,
  colback=orange!4!white,
  colframe=orange!80!black,
  arc=2mm, boxrule=0.8pt,
  left=2mm, right=2mm, top=1.5mm, bottom=1.5mm,
  title={\footnotesize\bfseries\color{white}\faSitemap\ SƠ ĐỒ TƯ DUY: MÔ HÌNH ĐOẠN THẲNG GẤP 4 LẦN}
]
\centering
\begin{tikzpicture}[
  seg/.style={draw=blue!80!black, fill=blue!15, rounded corners=1pt, minimum height=0.45cm, font=\scriptsize\bfseries}
]
  \node[anchor=east, font=\scriptsize\bfseries] at (0, 0.7) {Ngăn thứ nhất ($14$ cuốn):};
  \node[seg, minimum width=1.4cm] at (1.1, 0.7) {$14$};

  \node[anchor=east, font=\scriptsize\bfseries] at (0, 0) {Ngăn thứ hai (gấp $4$ lần):};
  \node[seg, minimum width=1.4cm, fill=teal!20, draw=teal!80!black] at (1.1, 0) {$14$};
  \node[seg, minimum width=1.4cm, fill=teal!20, draw=teal!80!black] at (2.6, 0) {$14$};
  \node[seg, minimum width=1.4cm, fill=teal!20, draw=teal!80!black] at (4.1, 0) {$14$};
  \node[seg, minimum width=1.4cm, fill=teal!20, draw=teal!80!black] at (5.6, 0) {$14$};

  \draw[thick, <->, purple!80!black] (0.4, -0.45) -- (6.3, -0.45) node[midway, below, font=\scriptsize\bfseries] {Ngăn thứ hai = $14 \times 4 = \mathbf{56}$ cuốn sách};
\end{tikzpicture}
\end{tcolorbox}

\vspace{-2.5mm}
% --- KHỐI 4: VÍ DỤ & MẸO TRÁNH BẪY ---
\noindent
\begin{minipage}[t]{0.49\textwidth}\vspace{0pt}
\begin{tcolorbox}[
  enhanced,
  colback=green!4!white,
  colframe=green!60!black,
  arc=2mm, boxrule=0.8pt,
  left=2mm, right=2mm, top=1.8mm, bottom=1.8mm,
  title={\footnotesize\bfseries\color{white}\faCheckCircle\ Ví dụ thực hành chuẩn mực}
]
{\footnotesize
\textbf{Bài toán:}\\
Một cửa hàng buổi sáng bán được $15\text{ kg}$ đường. Buổi chiều bán được số đường gấp $3$ lần buổi sáng. Hỏi buổi chiều cửa hàng bán được bao nhiêu ki-lô-gam đường?\\
\textbf{Bài giải:}\\
Buổi chiều cửa hàng bán được số ki-lô-gam đường là:
$$15 \times 3 = \mathbf{45}\text{ (kg)}$$
\textbf{Đáp số:} $45\text{ kg}$ đường.
}
\end{tcolorbox}
\end{minipage}
\hfill
\begin{minipage}[t]{0.49\textwidth}\vspace{0pt}
\begin{tcolorbox}[
  enhanced,
  colback=red!4!white,
  colframe=red!70!black,
  arc=2mm, boxrule=0.8pt,
  left=2mm, right=2mm, top=1.8mm, bottom=1.8mm,
  title={\footnotesize\bfseries\color{white}\faExclamationTriangle\ Mẹo tránh bẫy trắc nghiệm}
]
{\footnotesize
\begin{itemize}\setlength{\itemsep}{1.5pt}\setlength{\leftskip}{-2.5mm}
  \item \textbf{Bẫy hỏi cả hai buổi:}\\
  Đề bài hỏi "Cả hai buổi bán được bao nhiêu?", nhiều bạn chỉ tính số đường buổi chiều ($45\text{ kg}$) rồi dừng lại (thiếu bước cộng $15 + 45 = 60\text{ kg}$). Phải đọc kĩ câu hỏi!
  \item \textbf{Bẫy viết sai đơn vị:}\\
  Đơn vị phải là $(kg)$ chứ không được ghi là $(lần)$.
\end{itemize}
}
\end{tcolorbox}
\end{minipage}


\end{document}
